\setauthor{AC}
\section{Zentrale Verwaltung von Abonnements und Laufzeiten}
Mit der steigenden Zahl digitaler Streaming-Dienste wird auch die Verwaltung paralleler Abonnements komplexer. Nutzerinnen und Nutzer verwenden heute oft mehrere kostenpflichtige Angebote gleichzeitig, wodurch Monatskosten, Abrechnungszeitpunkte und inhaltliche Überschneidungen schwerer nachvollziehbar werden. Aktuelle Branchenanalysen zeigen, dass Haushalte regelmäßig mehrere Streaming-Abonnements parallel verwenden und dabei zunehmend auf Preis-Leistungs-Verhältnis, Nutzungsintensität und Kündigungsflexibilität achten \cite{deloitte_digital_media_2025}.

\section{Darstellung von Verfügbarkeiten von Filmen und Serien}
Neben der reinen Verwaltung von Abonnements ist für Anwender vor allem relevant, auf welcher Plattform ein bestimmter Film oder eine Serie verfügbar ist. Genau an dieser Stelle entsteht ein praktischer Mehrwert, wenn Inhalte und persönliche Abonnements gemeinsam betrachtet werden. Die verwendete TMDB API stellt dafür strukturierte Endpunkte für Filmdaten, Metadaten und Anbieterinformationen bereit und bildete damit eine fachliche Grundlage für die Umsetzung des Projekts \cite{tmdb_getting_started,tmdb_watch_providers}.

\section{Mehrwert einer zentralen Applikation für Nutzer}
Eine zentrale Applikation für Medien- und Abo-Management reduziert den organisatorischen Aufwand für Benutzer deutlich. Statt verschiedene Streaming-Plattformen einzeln zu überprüfen, lassen sich Suchfunktion, Verfügbarkeitsanzeige und eigene Abo-Daten in einer Anwendung zusammenführen. Gerade in einem Markt mit wachsender Anzahl an Diensten und steigender Preissensibilität wird ein solches zentrales Werkzeug für Benutzer zunehmend sinnvoll \cite{deloitte_digital_media_2025}.

\section{Datenschutz- und Sicherheitsanforderungen}
Sobald eine Anwendung persönliche Nutzungsdaten, Abonnements oder geschützte Funktionen bereitstellt, entstehen zusätzliche Anforderungen an Authentifizierung und Zugriffsschutz. Die Trennung zwischen öffentlichen Film- und Providerdaten einerseits und benutzerspezifischen Abo-Daten andererseits ist deshalb nicht nur eine technische, sondern auch eine fachliche Notwendigkeit. Für das Projekt war daraus abzuleiten, dass Benutzerkonten, Rollen und geschützte Requests von Beginn an als eigener Sicherheitsbereich betrachtet werden müssen.
